\begin{theorem}[Frobenius 局部积分叶定理]
设 $\mathcal D\subset TM$ 为光滑常秩 $k$ 分布，局部截面在 Lie 括号下封闭。每点附近存在坐标使\[\mathcal D=\operatorname{span}\{\partial_{y^1},\ldots,\partial_{y^k}\}.\]固定其余坐标得到局部积分叶。
\end{theorem}

\begin{proof}
目标是在分布对合的邻域上构造适配坐标使 $\mathcal D$ 由前 $k$ 个坐标方向张成。对秩 $k$ 归纳：$k=1$ 时，取非零局部截面 $X_1$，流盒定理给出坐标使
\[
X_1=\partial_{y^1},
\]
故积分曲线就是一维叶。

设结论对秩 $k-1$ 成立。取 $p$ 附近非零截面 $X_1\in\Gamma(\mathcal D)$，再次用流盒坐标令
\[
X_1=\partial_{y^1}.
\]
令横截面
\[
\Sigma:=\{y^1=0\}.
\]
通过对其余局部生成元减去适当的 $X_1$ 分量，可以选取
\[
Y_2,\ldots,Y_k\in\Gamma(\mathcal D)
\]
使它们在 $\Sigma$ 上切于 $\Sigma$。于是
\[
\mathcal E:=\mathcal D\cap T\Sigma
\]
是 $\Sigma$ 上秩 $k-1$ 分布。若 $Y,Z$ 是 $\mathcal E$ 的截面，则它们既切于 $\Sigma$ 又属于 $\mathcal D$。切于固定子流形的向量场在 Lie 括号下封闭，而 $\mathcal D$ 假设对合，因此
\[
[Y,Z]\in\Gamma(\mathcal D\cap T\Sigma)=\Gamma(\mathcal E).
\]
由归纳假设，缩小邻域后可在 $\Sigma$ 上取适配坐标
\[
(y^2,\ldots,y^k,y^{k+1},\ldots,y^n),
\]
使 $\mathcal E$ 由前 $k-1$ 个坐标方向张成；固定
$y^{k+1},\ldots,y^n$ 得到 $\mathcal E$ 的局部积分叶族 $L_c$。特别地，其中一片 $L_0$ 通过 $p$。

现在令 $\Phi_t$ 为 $X_1$ 的局部流。关键是证明它保持分布 $\mathcal D$。取局部标架 $X_1,\ldots,X_k$ 张成 $\mathcal D$。对合条件给出
\[
[X_1,X_a]=B^b{}_aX_b
\]
其中 $B^b{}_a$ 为光滑函数。沿 $X_1$ 的积分曲线，推前标架满足线性常微分方程；等价地，$(\Phi_{-t})_*X_a$ 的导数仍是这些向量的线性组合。因此初值位于 $\mathcal D$ 的向量在整个局部流中始终留在 $\mathcal D$，即
\[
(\Phi_t)_*\mathcal D=\mathcal D.
\]

对 $\Sigma$ 中每一片局部叶 $L_c$ 定义
\[
L_c^{\mathrm{sat}}
:=\{\Phi_t(q):q\in L_c,\ |t|<\varepsilon\}.
\]
其切空间由 $X_1$ 与 $(\Phi_t)_*T_qL_c$ 张成；前者属于 $\mathcal D$，后者由于分布被流保持也属于 $\mathcal D$。维数同为 $k$，故
\[
T_{\Phi_t(q)}L_c^{\mathrm{sat}}
=\mathcal D_{\Phi_t(q)}.
\]
因此这些饱和叶给出 $p$ 附近的局部叶层。把流参数作为 $y^1$，保留 $\Sigma$ 上的
$y^2,\ldots,y^n$，便得到局部坐标，其中
\[
\mathcal D
=\operatorname{span}\{\partial_{y^1},\ldots,\partial_{y^k}\}.
\]
这就得到 (iii)，并同时构造了局部积分叶。

$\mathbb R^3$ 中的分布 $\mathcal D=\operatorname{span}\{\partial_x,-y\partial_z+\partial_y\}$：$[\partial_x,-y\partial_z+\partial_y]=0\in\mathcal D$，对合，积分叶为 $z+\tfrac12y^2=\text{const}$ 的抛物柱面。

反例 $\mathcal D=\operatorname{span}\{\partial_x,x\partial_z+\partial_y\}$：$[\partial_x,x\partial_z+\partial_y]=\partial_z\notin\mathcal D$，不对合，不存在通过每点的二维积分曲面——这是接触几何的典型例子。

Frobenius 定理在控制论中给出可达性条件：向量场张成的分布对合当且仅当系统可沿积分叶任意移动。
\end{proof}

\begin{proposition}[余维一的 Pfaff 可积判据]
设 $\alpha$ 为处处非零光滑一形式。以下条件局部等价：$\ker\alpha$ 对合、$\alpha\wedge\dd\alpha=0$、存在 $\eta$ 使 $\dd\alpha=\eta\wedge\alpha$，以及存在非零 $\mu$ 和局部坐标 $u$ 使 $\alpha=\mu\dd u$。
\end{proposition}

\begin{proof}
余维一时分布 $\mathcal D=\ker\alpha$ 由单一一形式 $\alpha$ 确定。Frobenius 条件 $\dd\alpha=\eta\wedge\alpha$ 等价于 $\alpha\wedge\dd\alpha=0$，后者是可直接检验的代数条件。关键观察是：在 $n$ 维流形上，$\alpha\wedge\dd\alpha$ 是一个三形式，三形式的纤维维数为 $\binom n3$（仅 $n=3$ 时为一维），因此其消失条件只含 $\binom n3$ 个分量函数，远少于检验全部 Lie 括号。

先证 Frobenius $\Rightarrow$ Pfaff。若 $\dd\alpha=\eta\wedge\alpha$，则
\begin{equation*}
 \alpha\wedge\dd\alpha
 =\alpha\wedge\eta\wedge\alpha=0,
\end{equation*}
因为 $\alpha\wedge\alpha=0$（一形式与自身的楔积为零）。这一方向不需要 $\eta$ 的任何具体性质，只要它存在即可。

再证 Pfaff $\Rightarrow$ Frobenius。假设 $\alpha\wedge\dd\alpha=0$。局部取向量场 $N$ 使 $\alpha(N)=1$（归一化），并定义一形式
\begin{equation*}
 \eta:=-\iota_{N}\dd\alpha,
\end{equation*}
即 $\eta$ 是 $\dd\alpha$ 沿 $N$ 方向的缩并（带负号）。这里 $N$ 的存在性用到 $\alpha$ 处处非零，从而可在每点的余切空间中选一个对偶向量；局部光滑性由 $\alpha$ 光滑保证。

需在基向量对上逐一验证 $\dd\alpha=\eta\wedge\alpha$：取 $X,Y\in\ker\alpha$，并把 $\{X,Y,N\}$ 当作局部标架（在余维一情形下 $\ker\alpha$ 处处补上 $N$ 即给出 $TM$）：
\begin{itemize}
\item 在 $(X,Y)$ 上：由 $\alpha\wedge\dd\alpha=0$ 在 $(N,X,Y)$ 上取值，展开
\begin{equation*}
(\alpha\wedge\dd\alpha)(N,X,Y)
=\alpha(N)\dd\alpha(X,Y)-\alpha(X)\dd\alpha(N,Y)+\alpha(Y)\dd\alpha(N,X),
\end{equation*}
后两项因 $\alpha(X)=\alpha(Y)=0$ 消失，只剩 $\alpha(N)\dd\alpha(X,Y)=\dd\alpha(X,Y)=0$。而 $(\eta\wedge\alpha)(X,Y)=\eta(X)\alpha(Y)-\eta(Y)\alpha(X)=0$。两式相等。
\item 在 $(X,N)$ 上：$(\eta\wedge\alpha)(X,N)=\eta(X)\alpha(N)-\eta(N)\alpha(X)=\eta(X)=-\dd\alpha(N,X)=\dd\alpha(X,N)$。两式相等。这里最后一步用了 $\eta$ 的定义与二形式的反对称性。
\item 在 $(N,N)$ 上：两二形式在重复向量上取值均为零。
\end{itemize}
在所有基向量对上取值相同，故
\begin{equation*}
 \dd\alpha=\eta\wedge\alpha.
\end{equation*}

在余维一可积情形下，Frobenius 定理给出局部坐标 $(x^1,\ldots,x^{n-1},u)$ 使
\begin{equation*}
\alpha=\mu(x,u)\,\dd u,
\end{equation*}
其中 $\mu$ 是处处非零的光滑函数。叶状结构由 $u=\text{const}$ 给出，每片叶是 $n-1$ 维子流形。等价地，存在局部积分因子 $\mu^{-1}$ 使 $\mu^{-1}\alpha$ 成为恰当形式 $\dd u$；这正是经典 Pfaff 方程 $\alpha=0$ 求解的第一积分存在性。

接触结构的反例：在 $\RR^3$ 上取 $\alpha=\dd z-x\,\dd y$，则 $\dd\alpha=-\dd x\wedge\dd y$，于是
\begin{equation*}
\alpha\wedge\dd\alpha
=-(\dd z-x\,\dd y)\wedge\dd x\wedge\dd y
=-\dd z\wedge\dd x\wedge\dd y\neq0.
\end{equation*}
Pfaff 判据立即给出 $\ker\alpha$ 不可积，不存在二维积分曲面。这正是接触几何的出发点：$\alpha\wedge\dd\alpha\neq0$ 处处成立时 $\alpha$ 称为接触形式，$\ker\alpha$ 称为接触分布，它"最大限度地不可积"。

叶状结构的正例：在 $\RR^3$ 上取 $\alpha=\dd z-x\,\dd x$，则 $\dd\alpha=-\dd x\wedge\dd x=0$，从而 $\alpha\wedge\dd\alpha=0$。由于 $\alpha=\dd(z-\tfrac12x^2)$ 是恰当形式，等值面 $z-\tfrac12x^2=c$ 的切空间正是 $\ker\alpha$，因而它们构成所求的抛物柱面叶状结构。

于是\eqnrefs{eq:math-dg-frobenius-codim-one}与 Frobenius 条件局部等价。Pfaff 判据 $\alpha\wedge\dd\alpha=0$ 的优势在于不需要显式构造 $\eta$，只需检查一个三形式的消失；在余维一情形下它把可积性问题完全化为一个代数条件，这是接触几何与叶状结构理论分歧的精确分水岭。
\end{proof}

\begin{proposition}[接触分布的两步括号生成性]
设 $(M^{2m+1},\alpha)$ 为接触流形，$m\ge1$，$\mathcal H=\ker\alpha$。每点满足\[\mathcal H_p+[\mathcal H,\mathcal H]_p=T_pM,\qquad\alpha([X,Y])=-\dd\alpha(X,Y).\]
\end{proposition}

\begin{proof}
对水平向量场 $X,Y\in\Gamma(\mathcal H)$，有 $\alpha(X)=\alpha(Y)=0$。由\eqnrefs{eq:math-dg-frobenius-dalpha-bracket}，
\begin{equation}
\alpha([X,Y])=-\dd\alpha(X,Y).
\label{eq:math-dg-contact-bracket-vertical}
\end{equation}
固定一点 $p$ 和任意非零 $X_p\in\mathcal H_p$。由于接触条件等价于 $\dd\alpha|_{\mathcal H_p}$ 非退化，作为 $\mathcal H_p$ 上的辛形式，对每个非零 $X_p$ 必存在 $Y_p\in\mathcal H_p$ 使
\begin{equation*}
\dd\alpha_p(X_p,Y_p)\neq0.
\end{equation*}
具体地，取 $\mathcal H_p$ 的一组辛基 $\{e_i,f_j\}$ 使 $\dd\alpha_p(e_i,f_j)=\delta_{ij}$，将 $X_p$ 展开为 $X_p=\sum_i(a_i e_i+b_i f_i)$，则选 $Y_p=\sum_i(-b_i e_i+a_i f_i)$ 即给出
\begin{equation*}
\dd\alpha_p(X_p,Y_p)=\sum_i(a_i^2+b_i^2)=\|X_p\|^2>0.
\end{equation*}
把 $X_p,Y_p$ 延拓为局部水平场 $X,Y$（即 $\alpha(X)=\alpha(Y)=0$ 在邻域内成立），便由\eqnrefs{eq:math-dg-contact-bracket-vertical}得到
\begin{equation*}
\alpha_p([X,Y]_p)=-\dd\alpha_p(X_p,Y_p)\neq0.
\end{equation*}

设 $R_\alpha$ 为 Reeb 向量场，满足 $\alpha(R_\alpha)=1$、$\iota_{R_\alpha}\dd\alpha=0$。由 $\alpha_p([X,Y]_p)\neq0$，可写
\begin{equation*}
[X,Y]_p
=\alpha_p([X,Y]_p)\,R_\alpha(p)+H_p,
\qquad H_p\in\mathcal H_p,
\end{equation*}
所以 $[X,Y]_p$ 具有非零 Reeb 方向分量。又 $\mathcal H_p=\ker\alpha_p$ 维数为 $2m$，而 $T_pM$ 维数为 $2m+1$，故
\begin{equation*}
\dim\bigl(\mathcal H_p+[\mathcal H,\mathcal H]_p\bigr)
\ge\dim\mathcal H_p+1
=2m+1
=\dim T_pM,
\end{equation*}
反向不等式自动成立，即
\begin{equation}
\boxed{
\mathcal H_p+[\mathcal H,\mathcal H]_p=T_pM.
}
\label{eq:math-dg-contact-step-two-bracket-generating}
\end{equation}
接触分布因而是两步括号-生成分布：虽然单个允许速度永远限制在 $\mathcal H$ 中，但通过"先沿 $X$、再沿 $Y$、再反向返回"的小回路，可在二阶产生横向位移。
\end{proof}

\begin{proposition}[接触 Hamilton 场的辛化提升]
采用 $\alpha(X_H)=H$、$\iota_{X_H}\dd\alpha=(R_\alpha H)\alpha-\dd H$。在辛化 $\omega=\dd(\ee^s\alpha)$ 上按 $\iota_{X_{\widetilde H}}\omega=-\dd\widetilde H$ 定义场，则 $\widetilde H=\ee^sH$ 满足\[X_{\widetilde H}=X_H-(R_\alpha H)\partial_s.\]
\end{proposition}

\begin{proof}
在辛化上令
\[
\widetilde H(s,x)=\ee^sH(x).
\]
把 Hamiltonian 矢量场分解为
\[
X_{\widetilde H}=X_H+a\,\partial_s.
\]
利用\eqnrefs{eq:math-dg-contact-symplectization}，逐项计算
\[
\iota_{X_{\widetilde H}}\omega
=\ee^s\left[
a\alpha-H\dd s+\iota_{X_H}\dd\alpha
\right].
\]
而
\[
-\dd\widetilde H
=-\ee^s(H\dd s+\dd H).
\]
比较两式并代入\eqnrefs{eq:math-dg-contact-hamiltonian}，得到
\[
a=-R_\alpha H.
\]
所以
\begin{equation}
\boxed{
X_{\widetilde H}
=X_H-(R_\alpha H)\partial_s.
}
\label{eq:math-dg-contact-hamiltonian-lift}
\end{equation}

对常数 $\lambda>0$，记 $r_\lambda(s,x)=(s+\ln\lambda,x)$ 为沿 $s$ 方向的平移。由 $\widetilde H$ 的定义，
\begin{equation*}
\widetilde H\circ r_\lambda
=\ee^{s+\ln\lambda}H(x)
=\lambda\,\widetilde H(s,x),
\end{equation*}
即 $\widetilde H$ 是沿 $s$ 方向的 $1$ 次齐次函数。

对应地，辛 形式 $\omega=\dd(\ee^s\alpha)$ 在 $r_\lambda$ 下满足 $r_\lambda^*\omega=\lambda\,\omega$，因此 Hamiltonian 流保持这一齐次结构：若 $\phi_t$ 是 $\widetilde H$ 的流，则 $r_\lambda\circ\phi_t\circ r_{\lambda^{-1}}$ 也是 $\widetilde H$ 的流，由唯一性得 $r_\lambda\circ\phi_t=\phi_t\circ r_\lambda$。

Reeb 流作为特例：取 $H\equiv1$，则 $R_\alpha H=0$ 且 $X_H=R_\alpha$（由接触 Hamiltonian 的定义方程 $\iota_{X_H}\dd\alpha=(R_\alpha H)\alpha-\dd H=0$ 与 $\alpha(X_H)=H=1$）。此时
\begin{equation*}
X_{\widetilde H}=R_\alpha,
\end{equation*}
即 Reeb 流恰好提升为 辛化 上与 $s$ 坐标无关的 Hamiltonian 流。这给出 Reeb 流所有接触不变量的 辛 解释。

三维接触情形的显式公式：在标准接触 $\RR^3$ 上取 $\alpha=\dd z-\tfrac12(x\,\dd y-y\,\dd x)$，Reeb 场 $R_\alpha=\partial_z$。对 $H=H(x,y,z)$，接触 Hamiltonian 场为
\begin{equation*}
X_H
=\left(\frac{\partial H}{\partial y}+\frac{x}{2}\frac{\partial H}{\partial z}\right)\partial_x
-\left(\frac{\partial H}{\partial x}-\frac{y}{2}\frac{\partial H}{\partial z}\right)\partial_y
+\left(H-\frac{x}{2}\frac{\partial H}{\partial x}-\frac{y}{2}\frac{\partial H}{\partial y}\right)\partial_z,
\end{equation*}
提升到 $\RR^4_{(s,x,y,z)}$ 后沿 $s$ 方向的分量为 $-\partial_z H$。当 $H$ 与 $z$ 无关时提升为齐次 Hamiltonian 系统，这正是经典 Hamilton 力学在接触框架下的对应。

接触 Hamilton 系统也能编码摩擦，但必须保持同一符号约定。
在 $\alpha=\dd z-p\dd q$ 上，本章采用 $\alpha(X_H)=H$，因此
\[
\dot q=-H_p,\qquad\dot p=H_q+pH_z,
\qquad\dot z=H-pH_p,
\qquad\dot H=(R_\alpha H)H.
\]
若要得到 $\ddot q+\gamma\dot q+\omega^2q=0$、$\gamma>0$，应取
\[
H=-\left(\tfrac12p^2+\tfrac12\omega^2q^2+\gamma z\right).
\]
于是 $\dot q=p$、$\dot p=-\omega^2q-\gamma p$。
机械能 $E=\tfrac12p^2+\tfrac12\omega^2q^2$ 满足
\[
\dot E=-\gamma p^2\le0,\qquad\dot H=-\gamma H.
\]
若使用原来的正号 $H=E+\gamma z$，在当前约定下得到的是反摩擦，
不能把 $\dot H=\gamma H$ 解释成摩擦能量耗散。
辛化后 $\widetilde H=\ee^sH$ 保守，且 $\dot s=-R_\alpha H=\gamma$；
额外坐标抵消接触能量的指数因子，保留 $\widetilde H$。

接触 Hamiltonian 动力学 因而可以看成一类对 $s$ 方向齐次的 辛 Hamiltonian 动力学；Reeb 流 对应 $H\equiv1$ 的特例，而一般情形下 Reeb 分量 $-R_\alpha H$ 正好度量系统偏离保守的程度。
\end{proof}

\begin{theorem}[Ehresmann 局部平凡化定理]
设 $F:M\to B$ 为光滑、满射、真的下沉映射。每个 $b_0$ 附近存在邻域 $U$ 与微分同胚\[\Phi:U\times F^{-1}(b_0)\to F^{-1}(U),\qquad F\circ\Phi(b,x)=b.\]
\end{theorem}

\begin{proof}
在 $b_0$ 附近选坐标球 $U$ 使其闭包紧，并对每个 $b\in U$ 取从 $b_0$ 指向 $b$ 的径向路径 $\gamma_b:[0,1]\to U$，满足 $\gamma_b(0)=b_0$、$\gamma_b(1)=b$。把 $\dot\gamma_b(t)$ 水平提升到 $M$，从任意 $x\in F^{-1}(b_0)$ 出发解 ODE
\[
\dot{\widetilde\gamma}_{b,x}(t)
=\dot\gamma_b(t)^H,
\qquad
\widetilde\gamma_{b,x}(0)=x.
\]
沿提升曲线计算 $F$ 的变化率，由链式法则与水平提升的定义，
\begin{align*}
\frac{\dd}{\dd t}F(\widetilde\gamma_{b,x}(t))
&=\dd F_{\widetilde\gamma_{b,x}(t)}\bigl(\dot{\widetilde\gamma}_{b,x}(t)\bigr)\\
&=\dd F_{\widetilde\gamma_{b,x}(t)}\bigl(\dot\gamma_b(t)^H\bigr)
=\dot\gamma_b(t),
\end{align*}
最后一步用到 $\dd F|_{\mathcal H}$ 是恒等提升。积分两端并代入初值 $F(\widetilde\gamma_{b,x}(0))=F(x)=b_0=\gamma_b(0)$，得到
\[
F(\widetilde\gamma_{b,x}(t))=\gamma_b(t),
\]
因此提升曲线始终停留在对应纤维之上。因为 $F$ 真，紧小球 $\overline U$ 的原像 $F^{-1}(\overline U)$ 紧，水平 ODE 的解不会在有限时间逃向无穷远；缩小 $U$ 后，所有这些 ODE 在 $t\in[0,1]$ 上统一存在。定义
\[
\Phi(b,x):=\widetilde\gamma_{b,x}(1).
\]
由 ODE 对初值 $x$ 与参数 $b$ 的光滑依赖，$\Phi$ 光滑；且由上式 $F(\Phi(b,x))=\gamma_b(1)=b$。构造逆映射时，对 $(y,b)\in F^{-1}(U)$ 取从 $b$ 沿径向返回 $b_0$ 的路径 $\gamma_b(1-t)$，从 $y$ 出发水平提升至 $t=1$，得到 $F^{-1}(b_0)$ 中的一点，与 $b$ 组合即给出 $\Phi^{-1}$，它同样由 ODE 光滑依赖保证光滑。因此
\[
\Phi:U\times F^{-1}(b_0)\to F^{-1}(U)
\]
是微分同胚，并满足 $F\circ\Phi(b,x)=b$，这正得到\eqnrefs{eq:math-dg-ehresmann-local-trivialization}。
\end{proof}

\begin{proposition}[主联络曲率与水平可积性]
对主联络 $\omega$，令 $\Omega=\dd\omega+\frac12[\omega\wedge\omega]$、$\mathcal H=\ker\omega$。则\[\Omega(X^H,Y^H)=-\omega([X^H,Y^H]),\qquad\Omega=0\Longleftrightarrow\mathcal H\text{ 对合}.\]全局水平截面还要求和乐平凡。
\end{proposition}

\begin{proof}
对任意 $\mathfrak g$-值一形式 $\omega$ 与任意向量场 $X,Y$，外微分的 不变 公式为
\[
 \dd\omega(X,Y)
 =X\{\omega(Y)\}-Y\{\omega(X)\}-\omega([X,Y]).
\]
令 $X=X^H$、$Y=Y^H$ 为水平提升。由水平分布定义 $\mathcal H=\ker\omega$，有 $\omega(X^H)=\omega(Y^H)=0$，前两项均为零，剩下
\[
 \dd\omega(X^H,Y^H)=-\omega([X^H,Y^H]).
\]
再算楔括号项。由定义
\[
 [\omega\wedge\omega](X^H,Y^H)
 =[\omega(X^H),\omega(Y^H)]-[\omega(Y^H),\omega(X^H)],
\]
每个因子均为零，故
\[
 [\omega\wedge\omega](X^H,Y^H)=0.
\]
把两式代入曲率定义 $\Omega=\dd\omega+\tfrac12[\omega\wedge\omega]$，得到
\begin{align*}
\Omega(X^H,Y^H)
&=\dd\omega(X^H,Y^H)+\tfrac12[\omega\wedge\omega](X^H,Y^H)\\
&=-\omega([X^H,Y^H])+0,
\end{align*}
即\eqnrefs{eq:math-dg-curvature-nonintegrability}。

若 $\Omega=0$，则对一切水平对有 $\omega([X^H,Y^H])=0$，即 $[X^H,Y^H]\in\ker\omega=\mathcal H$，括号仍水平，Frobenius 定理给出局部水平叶。

反之若水平分布 对合的，则 $[X^H,Y^H]\in\mathcal H$，故 $\omega([X^H,Y^H])=0$，由上式曲率在全部水平对上消失；主联络曲率又满足 horizontality，即它被水平向量场完全决定，因此 $\Omega=0$ 处处成立，即\eqnrefs{eq:math-dg-horizontal-frobenius}。

第一 Bianchi 恒等式 $D\Omega=0$ 可视为水平分布的"高阶可积性"约束：若 $\Omega=0$ 给出局部叶，则 $D\Omega=0$ 自动满足；若 $\Omega\neq0$，Bianchi 限制曲率沿水平方向变化的可能方式。在联络语言中 Bianchi 恒等式是 $D\Omega=0$。
协变外微分一般并不满足 $D^2=0$：在伴随值形式上
$D^2\eta=[\Omega\wedge\eta]$，其平方恰由曲率控制。
黎曼第一 Bianchi 还使用无挠条件，第二 Bianchi 则是协变曲率恒等式。

当结构群 $G$ 为 Abelian（如 $U(1)$ 电磁联络），$[\omega\wedge\omega]=0$，曲率简化为 $\Omega=\dd\omega$。

此时 Frobenius 障碍就是 $\dd\omega=0$，即 $\omega$ 闭；局部叶存在当且仅当 $\omega$ 局部恰当，这正是电磁势的规范不变性 $\omega\mapsto\omega+\dd f$ 不改变曲率。

非 Abelian 情形下 $[\omega\wedge\omega]$ 项给出非交换性贡献，使曲率不再是线性算子，Yang--Mills 理论的非线性性正源于此。

对 $U(1)$ 联络在圆 $S^1$ 上的平行移动，沿闭路 $\gamma$ 的和乐为
\begin{equation*}
\operatorname{Hol}(\gamma)=\exp\left(i\oint_\gamma\omega\right)
=\exp\left(i\int_D\Omega\right),
\end{equation*}
最后一步要求闭路在相应平凡化内确实为某曲面 $D$ 的边界，才可用 Stokes 定理；圆周上的一般非可缩闭路没有这样的 $D$。曲率 $\Omega$ 因此是和乐的"密度"：把闭路包围的曲率积分即得和乐角度。非 Abelian 情形下此式变为非交换 Wilson 圈
\begin{equation*}
W(\gamma)=\mathcal P\exp\left(\oint_\gamma\omega\right),
\end{equation*}
其中 $\mathcal P$ 表示路径序，反映非交换性使局部贡献不能简单交换。

但 $\Omega=0$ 仍只是局部可积性。固定 $p_0\in P$ 后，沿基流形路径作水平提升定义平行移动。若基空间存在非平凡闭路，平坦联络仍可能具有非平凡和乐表示
\[
 \operatorname{Hol}:\pi_1(B,b_0)\to G.
\]
只有当这些闭路的和乐全部平凡时，从 $p_0$ 出发的水平叶才会在每个基点上选出唯一一点，从而成为全局水平截面。若 $B$ 单连通，则 $\pi_1(B)=0$，平坦性使平行移动与路径无关，于是该局部 Frobenius 结构可以沿整个连通基底拼成全局水平截面。由此可见：曲率控制局部可积性，和乐 控制全局拼接；前者是微分条件，后者是拓扑条件，两者共同决定联络的可积性结构。

曲率与和乐的分离对应特征类的两层结构：第一 Chern 类 $c_1\in H^2(B,\ZZ)$ 通过 de Rham 表示由曲率积分给出（局部信息），但其整体挠部分由 和乐 决定（全局信息）。Chern--Weil 理论把曲率不变多项式映射到上同调类，给出 $c_k$、$p_k$、Euler 类等的微分代表；积分特征类的挠部分不能由实曲率形式检出，需整体转移函数或微分特征类等额外数据。$\hat A$-亏格和通常的数值指标同样不能检出一般挠类，不应把它们当作完整恢复挠信息的工具。
\end{proof}

\begin{proposition}[无穷小矩形和乐的曲率系数]
局部取 $D=\dd+A$，以列向量满足 $\dot U=-A(\dot\gamma)U$。沿正向 $X,Y,-X,-Y$ 的坐标小矩形，\[U_{\partial\square}=I-\varepsilon^2F(X,Y)+O(\varepsilon^3),\qquad F=\dd A+A\wedge A.\]
\end{proposition}

\begin{proof}
取局部联络形式 $A=A_i\,\dd x^i$，沿 $X$、$Y$ 方向各走长度 $\varepsilon$ 的小矩形。平行移动在每条边上由 $\dot U=-A(\dot\gamma)U$ 给出。
为完整保留二阶项，所有 $A$ 及其导数都在起点计算，四条边分别为
\[
\begin{aligned}
U_1&=I-\varepsilon A_X+\tfrac{\varepsilon^2}{2}(A_X^2-\partial_XA_X),\\
U_2&=I-\varepsilon A_Y+\varepsilon^2(\tfrac12A_Y^2-\partial_XA_Y-\tfrac12\partial_YA_Y),\\
U_3&=I+\varepsilon A_X+\varepsilon^2(\tfrac12A_X^2+\tfrac12\partial_XA_X+\partial_YA_X),\\
U_4&=I+\varepsilon A_Y+\tfrac{\varepsilon^2}{2}(A_Y^2+\partial_YA_Y).
\end{aligned}
\]
每式余项均为 $O(\varepsilon^3)$；对列向量，完整乘积为
$U_{\partial\square}=U_4U_3U_2U_1$。
原来只保留每边的一阶项再声称精确到二阶，会遗留错误的 $A_X^2,A_Y^2$ 项；
这些同边项必须与上面的二阶指数项同时计算。
一次项全部消失，二次项整理为
\[
U_{\partial\square}
=I-\varepsilon^2
\left(
\partial_XA_Y-
\partial_YA_X+[A_X,A_Y]
\right)
+O(\varepsilon^3).
\]
括号正是曲率分量
\[
F(X,Y)
=(\dd A+A\wedge A)(X,Y).
\]
因此曲率不是平行移动误差的比喻，而是无穷小闭路和乐的首个非平凡系数。

把四条边的乘积按 $\varepsilon$ 阶数展开。零阶项为 $I$。一阶项为
\begin{equation*}
-\varepsilon A_X-\varepsilon A_Y+\varepsilon A_X+\varepsilon A_Y=0,
\end{equation*}
正反向贡献相消，反映沿闭路回到起点的最低阶必然平凡。二阶项分三类：
\begin{itemize}
\item 同向导数项：$-\varepsilon^2\partial_XA_Y$ 与 $+\varepsilon^2\partial_YA_X$，来自沿一边走时联络值在另一边方向的变化；
\item 交换项：$-\varepsilon^2[A_X,A_Y]$，来自非交换 Lie 代数的乘积顺序差异；
\item 同边往返项：如 $(-\varepsilon A_X)(\varepsilon A_X)=-\varepsilon^2A_X^2$，但这些项在同一条边往返时已被路径序局部解吸收，不出现在最终公式中。
\end{itemize}
三类贡献合并即得曲率分量。

当 $G$ 为 Abelian（如 $U(1)$），$[A_X,A_Y]=0$，曲率简化为 $F_{XY}=\partial_XA_Y-\partial_YA_X$，这正是电磁场强 $F_{\mu\nu}=\partial_\mu A_\nu-\partial_\nu A_\mu$。沿小回路的和乐为
\begin{equation*}
U_{\partial\square}=\exp\left(-i\varepsilon^2F_{XY}\right)+O(\varepsilon^3),
\end{equation*}
对应回路包围的磁通；只有进一步考虑随时间变化的回路或场，才联系到 Faraday 感应定律。

对 $SU(N)$ 联络，交换项 $[A_X,A_Y]$ 给出非线性自相互作用，是 Yang--Mills 方程
\begin{equation*}
D^\mu F_{\mu\nu}=0
\end{equation*}
非线性的来源。物理上这对应胶子自相互作用：与电磁场光子不带电荷不同，胶子本身带色荷，因此曲率含 $A\wedge A$ 项。

把无穷小回路公式积分到有限回路，得到和乐沿任意闭路 $\gamma$ 的表达式
\begin{equation*}
\operatorname{Hol}(\gamma)=\mathcal P\exp\left(-\oint_\gamma A\right),
\end{equation*}
对边界为 $\gamma$ 的曲面 $\Sigma$ 应用非交换 Stokes 定理，得
\begin{equation*}
\operatorname{Hol}(\gamma)=\mathcal P_\Sigma\exp\left(-\int_\Sigma \widehat F\right),\qquad \widehat F(x)=U_x^{-1}F(x)U_x,
\end{equation*}
这里先选择从同一基点到 $x$ 的扫描路径，以 $U_x$ 把各点曲率运输回同一纤维，再采用与扫描相容的曲面积序 $\mathcal P_\Sigma$。

不能直接把不同纤维中的 $F(x)$ 相乘。

把 $\Sigma$ 分成 $N^2$ 个面积为 $1/N^2$ 的小方格并取 $N\to\infty$，每个小方格贡献 $I-\frac1{N^2}F+O(N^{-3})$，乘积收敛到上式。

该式适用于边界为 $\gamma$ 的这类曲面与选定扫描。

Ambrose--Singer 定理则说受限和乐的 Lie 代数由平行运输回基点的曲率值生成；非可缩闭路上的整体和乐仍需独立拓扑数据。

在量子力学中，参数空间上的绝热演化给出 $U(1)$ 联络（Berry 联络）$\mathcal A=i\langle\psi|\dd\psi\rangle$，曲率 $\mathcal F=\dd\mathcal A$ 是 Berry 曲率。沿参数空间闭路演化一周后波函数获得相位 $\exp(i\oint\mathcal A)$，这正是无穷小回路公式在 $U(1)$ 情形的应用。Aharonov--Bohm 效应中电子绕磁通量管的相位差则是非平凡 和乐 的典型例子：联络局部平坦（磁场为零区域）但全局 和乐 非平凡，体现曲率与 和乐 的全局-局部分离。
\end{proof}
